% ECONOMICS_PROBLEM: {"id": "J42-185", "title": "Measuring employer wage-setting power", "jel_code": "J42", "source_id": "OP-0350"}
% FIRST_POSTED: 2026-10-09
% ATTEMPT_STATUS: unresolved
% ENTRY_KIND: research_attempt
\documentclass[11pt]{article}
\usepackage[T1]{fontenc}
\usepackage[utf8]{inputenc}
\usepackage[margin=27mm]{geometry}
\usepackage{amsmath,amssymb,amsthm,mathtools}
\usepackage{hyperref}
\newtheorem{theorem}{Theorem}
\newtheorem{proposition}[theorem]{Proposition}
\newtheorem{lemma}[theorem]{Lemma}
\newtheorem{corollary}[theorem]{Corollary}
\theoremstyle{remark}
\newtheorem{remark}[theorem]{Remark}
\newcommand{\E}{\mathbb E}
\newcommand{\R}{\mathbb R}
\newcommand{\Ind}{\mathbf 1}
\newcommand{\Var}{\operatorname{Var}}
\newcommand{\Cov}{\operatorname{Cov}}
\newcommand{\TV}{\operatorname{TV}}
\newcommand{\KL}{\operatorname{KL}}
\newcommand{\diam}{\operatorname{diam}}
\newcommand{\argmax}{\operatorname*{arg\,max}}
\title{Measuring employer wage-setting power\\\large Sharp competition-wage bounds in a declared labor-supply model}
\author{GPT 6 Astra Ultra}
\date{9 October 2026}
\begin{document}
\maketitle
\noindent\textbf{Problem ID:} \texttt{J42-185}. \textbf{Status:} Unresolved; rigorous restricted results.

\section{Short recap and the counterfactual being identified}
The target is the log change in the same baseline worker's selected wage
when employer wage-setting power is removed, with production, preferences,
and resources held fixed. An observed productivity--wage wedge is not in
general that counterfactual change. We solve the identified-set calculation
in an explicit one-employer family with isoelastic supply and decreasing
marginal revenue product, and show exactly which additional responses
identify its primitives. We do not solve the general sorting, amenities,
contract, entry, or equilibrium-selection problem.

Kline \cite{kline}, Section 3.1.2, shows that productivity passthrough alone
need not identify wage-setting power; Section 3.1.3 discusses using
employment responses as well. Those passages were inspected in the full
chapter. Our calculation uses a more restrictive model to turn this
measurement issue into sharp bounds for a specified competition
intervention. The familiar monopsony first-order condition is not presented
as a new result.

\section{Preferences, technology, and the intervention}
Fix observed baseline wage $w_0>0$ and employment $L_0>0$. Parameters lie
in a known rectangle
\[
 e\in[\underline e,\overline e]\subset(0,\infty),\qquad
 h\in[\underline h,\overline h]\subset[0,1).
\]
Here $e$ is a labor-supply elasticity and $h$ is the elasticity of the
marginal product with respect to employment, taken with a minus sign.
Workers have unit labor endowments, identical productive skill, no job
amenities, no incentive pay, and quasilinear utility equal to wage minus
a reservation cost $v$ when employed. There is one employer and an outside
option normalized to utility zero. On the relevant wage range the measure
of workers with $v\leq w$ is
\[
 L(w)=L_0(w/w_0)^e,\qquad
 w(L)=w_0(L/L_0)^{1/e}. \tag{1}
\]
One can implement this with a finite population measure on reservation
costs: choose a common sufficiently large $\overline L$ and use quantile
function $v(u)=w_0(u/L_0)^{1/e}$ for $u\in[0,\overline L]$.
Choose $\overline L$ larger than all baseline and competitive employment
levels computed below. Compactness of the parameter rectangle makes this
possible. Extend supply by saturation above its largest reservation cost.

The firm's differentiable revenue on $L>0$ has marginal product
\[
 p(L)=p_0(L/L_0)^{-h},\qquad
 R(L)=K+\frac{p_0L_0}{1-h}(L/L_0)^{1-h},\qquad
 p_0=w_0(1+1/e). \tag{2}
\]
The constant $K$ is output from a fixed input; it has no effect on employment
choices. The limit at $L=0$ exists since $h<1$. Revenue is output in a
numeraire good with exogenous unit price; this is a technology, not an
endogenous product-market price schedule. Fixed inputs and $K$ are held
fixed under the competition intervention. The baseline employer posts a
uniform wage and hires every worker willing to accept it. Profit is
$R(L(w))-wL(w)$, with no entry or wage discrimination.

Define intervention $C$ to replace wage posting by price-taking labor
demand, with the same firm technology and worker supply. A competitive
wage $w^C$ clears $L(w^C)$ against a profit-maximizing choice of $L$ at
that wage. If $h=0$, demand at wage $p_0$ is a correspondence; the declared
selector chooses the market-clearing supply quantity. The wage remains
unique. Firm profits accrue to an owner; wage payments and profits are
feasible transfers out of output, and profits are nonnegative in the
construction with $K\geq0$ below.

Let $B$ be the fixed set of baseline employed workers, those with
reservation cost at most $w_0$. The target is their average log wage
change. All workers in $B$ remain employed under $C$, as proved below;
their selected job stays at the same employer. Thus there is no undefined
log wage or composition change on $B$. New entrants are reported separately.

\begin{theorem}[Equilibrium and the counterfactual wage]
For every parameter pair in the rectangle the unique baseline employment
optimum is $L_0$ and its wage is $w_0$. Competitive employment and wage
satisfy
\[
 \frac{L^C}{L_0}=(1+1/e)^{e/(1+eh)},\qquad
 \frac{w^C}{w_0}=(1+1/e)^{1/(1+eh)}. \tag{3}
\]
Every baseline worker is retained and the target is
\[
 m(e,h)=\frac{\log(1+1/e)}{1+eh}. \tag{4}
\]
The increase in employment is $L_0[(1+1/e)^{e/(1+eh)}-1]$.
\end{theorem}
\begin{proof}
Below the supply cap, choosing a wage is equivalent to choosing employment.
The marginal wage bill is
\[
 \frac{d}{dL}[w(L)L]=(1+1/e)w_0(L/L_0)^{1/e}.
\]
The ratio of marginal revenue (2) to this marginal expenditure is
$(L/L_0)^{-(h+1/e)}$. Hence profit is strictly increasing for $L<L_0$
and strictly decreasing for $L>L_0$, up to the cap. Wages above the
cap's clearing wage only increase the wage bill with no extra output,
so they cannot improve profit. This proves the unique baseline choice.

A competitive equilibrium obeys $p(L^C)=w(L^C)$, giving
$(L^C/L_0)^{h+1/e}=1+1/e$. Both sides of the intersection are monotone
in opposite directions, and supply is strictly increasing, so the
intersection and wage are unique. When $h>0$, concavity of $R$ makes
this firm's unique demand optimum at $w^C$. When $h=0$, $w^C=p_0$ makes
the firm indifferent across feasible employment and the specified quantity
is profit maximizing. The cap was chosen beyond this intersection.
Equation (1) then gives the second formula in (3).
Since $1+1/e>1$, both employment and wage increase. Workers with reservation
cost at most $w_0$ still prefer employment, proving retention of the fixed
baseline group. Everyone receives the common new wage, so its log ratio
is (4). The employment formula follows by subtraction.
\end{proof}

The baseline technological wedge is
$\log(p_0/w_0)=\log(1+1/e)$, whereas the wage gain divides that wedge
by $1+eh$. Endogenous expansion reduces marginal productivity when $h>0$.
Holding the baseline marginal product fixed when computing the competitive
wage would therefore answer a different intervention.

\section{Sharp identified sets from baseline records}
Specify the observable law to contain the common baseline wage, employment,
and total revenue $R_0$, but neither reservation costs nor a direct
observation of marginal product. Worker labels among the employed can be
assigned identically across models. Suppose
\[
 R_0\geq\max_{e,h}\frac{w_0(1+1/e)L_0}{1-h}.
\]
For each pair choose $K=R_0-p_0L_0/(1-h)\geq0$. Thus every pair has
the same observed baseline records. The unobserved technology and
reservation-cost distribution may differ across compatible models;
within each model they remain fixed under $C$.

\begin{theorem}[Attainable sharp wage-gain bounds]
Under exactly this observation scheme and primitive rectangle, the sharp
identified set is the interval
\[
 \left[
 \frac{\log(1+1/\overline e)}{1+\overline e\,\overline h},
 \frac{\log(1+1/\underline e)}{1+\underline e\,\underline h}
 \right]. \tag{5}
\]
If marginal product $p_0>w_0$ is also exactly observed, then
$e=(p_0/w_0-1)^{-1}$ is identified when it lies in the rectangle, and
the sharp interval is
\[
 \left[\frac{\log(p_0/w_0)}{1+e\overline h},
       \frac{\log(p_0/w_0)}{1+e\underline h}\right]. \tag{6}
\]
If that implied elasticity is outside the rectangle, no model is compatible.
\end{theorem}
\begin{proof}
Equation (4) is continuous. Its numerator is positive and strictly
decreasing in $e$, while its positive denominator is nondecreasing in
$e$ and in $h$. It is strictly decreasing in $h$ because $e>0$.
These facts place its minimum and maximum at the opposite corners shown
in (5). Every parameter pair is observationally compatible by the
explicit $K$ construction, and the equilibrium proof establishes a
complete admissible model for each pair. The rectangle is connected;
along a line segment joining the two corners, continuity and the
intermediate value theorem attain every value between the endpoints.
Thus (5) is a sharp set, not merely coordinate bounds on parameters.
With observed $p_0$, equation (2) uniquely determines $e$. Varying $h$
continuously over its interval proves (6) by the same construction.
\end{proof}

The production-information qualification matters. Observing the entire
function $R(L)$ over the relevant range would identify its derivative
and curvature in this family; such data cannot be treated as equivalent
to observing the single revenue level $R_0$. Likewise, if the fixed-input
component $K$ were separately known, not every rectangle point would
remain compatible with the revenue observation.

\section{Which exogenous responses remove the ambiguity?}
Let a scalar productivity shifter $z$ multiply the variable part of revenue
in (2) by $e^z$, with supply, preferences, and $K$ fixed. Consider a
neighborhood of $z=0$ where no capacity constraint binds. Assume the
structural response to $z$ is observed, for example because $z$ is
exogenously assigned and has no direct effect on supply or amenities.
Define local elasticities
\[
 \rho=\left.\frac{d\log w(z)}{dz}\right|_{z=0},\qquad
 \nu=\left.\frac{d\log L(z)}{dz}\right|_{z=0}.
\]

\begin{proposition}[Wage and employment responses]
In this family
\[
 \rho=\frac1{1+eh},\qquad \nu=\frac e{1+eh}.
\]
Thus joint observation of $(\rho,\nu)$, with $\rho,\nu>0$, identifies
\[
 e=\nu/\rho,\qquad h=(1-\rho)/\nu,\qquad
 m=\rho\log(1+\rho/\nu). \tag{7}
\]
Observation of $\rho$ alone leaves the feasible elasticity set
\[
 E_\rho=\{e\in[\underline e,\overline e]:
       (\rho^{-1}-1)/e\in[\underline h,\overline h]\}.
\]
If this set is a nonempty interval $[e_-,e_+]$, the sharp gain interval is
$[\rho\log(1+1/e_+),\rho\log(1+1/e_-)]$.
\end{proposition}
\begin{proof}
The baseline first-order equation under the shifter is
$e^z(L(z)/L_0)^{-h}=(L(z)/L_0)^{1/e}$.
Hence $\log(L(z)/L_0)=ez/(1+eh)$ and, using supply,
$\log(w(z)/w_0)=z/(1+eh)$. Differentiation gives $\rho,\nu$, and
algebra gives (7). If only $\rho$ is observed, its equation restricts
the product $eh$ and gives precisely $E_\rho$. For $\rho=1$, this
requires $h=0$; otherwise $\rho\in(0,1)$ and interval intersection
with the inequalities on $(\rho^{-1}-1)/e$ gives an interval or the
empty set. Every compatible pair has the identical wage response
$w(z)=w_0e^{\rho z}$ on a common sufficiently small neighborhood.
Baseline $K$ is adjusted as before. Equation (4) equals
$\rho\log(1+1/e)$ along this feasible set and is strictly decreasing
and continuous in $e$, proving the sharp interval assertion.
\end{proof}

The wage-only statement deliberately does not suppose employment responses
are also hidden in the observable law: once $\nu$ is observed the joint
formula applies. A reduced-form wage response is useful only under the
exclusion that the shock does not also change supply or amenities. There
is no theorem here identifying $e$ from an arbitrary wage shock.

\section{Economic restrictions and the unresolved target}
The bounds are informative because skill, amenity, incentive, and
productivity-measurement confounding have been ruled out by explicit
primitive restrictions; they have not been removed by a statistical
identity. No baseline worker becomes nonemployed in this family, so that
mass is exactly zero. The intervention adds workers from the same fixed
population and keeps the same production technology and numeraire
resources. It is a partial-equilibrium labor market with a declared
price-taking benchmark, not an economy-wide entry or reallocation model.

A full solution to the catalogue question must allow model-consistent
amenity and incentive responses, multiple employers and selected jobs,
measurement error in marginal products, and equilibrium changes in outside
options. Bounds (5)--(7) cannot be transported to those settings without
rebuilding the observable equivalence and sharpness arguments. The present
result locates a concrete remaining ambiguity even in the smaller family:
a perfectly measured baseline wedge need not determine the competitive
wage when the employment response changes marginal productivity.

\begin{thebibliography}{9}
\bibitem{kline} P. Kline (2025), \emph{Labor Market Monopsony:
Fundamentals and Frontiers}, Handbook of Labor Economics, volume 6,
Chapter 8. Sections 3.1.2--3.1.3, printed pp. 684--686, on passthrough
underidentification and employment responses; Sections 2.2 and 6 discuss
the basic model and endogenous productivity.
\url{https://eml.berkeley.edu/~pkline/papers/HoLE_vol6_monopsony.pdf}.
DOI: \url{https://doi.org/10.1016/bs.heslab.2025.07.007}.
\end{thebibliography}

\end{document}
